\scanpage{245}
\mspage{423}

\margin{one should \unsure{give its} \uline{formulas} in a precise \unsure{way}, \unsure{take} \illegible{2}!!!}

Now according to the general principles (not yet established, to tell the truth, in the profinite setting) the commutant of $\rho$ in the group in question is reduced to the group generated by $\rho$, $\{1,\rho,\rho^2\}$, so that the relation \unsure{(*)} takes the form
\[
\tau_\infty^{\prime\alpha_3}\widehat{\varphi}(f_\infty)\alpha=\rho^j
\qquad\text{for }j\in\mathbb Z/3\mathbb Z\ \unsure{\text{suitable}}
\quad\text{(well determined)}.
\]
We make it explicit by \uline{using} the value of $f_\infty$
\[
f_\infty^{-1}=g\rho(g)\rho^2(f_0^{-1})
=g\rho(g)=\rho^2(g)^{-1}=\rho^{-1}(\xi^{-1}),
\]
\inserted{N.B. here $\left\{\begin{gathered}f_0=f=1\\g=\xi=\alpha\rho(\alpha)^{-1}\end{gathered}\right.$}
thus the relation \deleted{\illegible{2}} becomes
\[
\tau_\infty^{\prime\alpha_3}\widehat{\varphi}\bigl(\rho^{-1}(\xi)\bigr)\alpha=\rho^j
\]
i.e., \deleted{by applying $\tau_\infty^{\prime\alpha_3}$ \unsure{to the}}
\[
\widehat{\varphi}\bigl(\rho^{-1}(\xi)\bigr)
=\underbrace{\tau_\infty^{\prime\alpha_3}\rho^j}_{\underbrace{\tau_\infty^{\prime\alpha_3}(\rho^j)}\tau_\infty^{\prime\alpha_3}}\alpha^{-1}
\tag{\deleted{\unsure{44|43}}}
\]
or also
\deleted{\[
\widehat{\varphi}\bigl(\rho^{-1}(\xi^{-1})\bigr)=\alpha\rho^{-j}\tau_\infty^{\prime\alpha_3}.
\tag{\unsure{44 bis|43 bis}}
\]}
since $\tau'_\infty(\rho)=\rho^{-1}$ thus
$\tau_\infty^{\prime\alpha_3}(\rho^j)=\rho^{(-1)^{\alpha_3}j}=\rho^{j'}$
\[
\widehat{\varphi}\bigl(\rho^{-1}(\xi)\bigr)=\rho^{j'}\tau_\infty^{\prime\alpha_3}\alpha^{-1}
\tag{44}
\]
where $\alpha_3=\alpha_3(\mu)\in\mathbb Z/2\mathbb Z$ defined by $(-1)^{\alpha_3}\equiv\mu\ (3)$ and $j'\in\mathbb Z/3\mathbb Z$ (well determined).

We shall transform this relation a little further, by writing
\[
\tau_\infty^{\prime\alpha_3}=(\sigma_\infty\tau_\infty)^{\alpha_3}=\sigma_\infty^{\alpha_3}\tau_\infty^{\alpha_3}
\qquad\text{whence}
\]
\[
\tau_\infty^{\prime\alpha_3}\alpha^{-1}=\sigma_\infty^{\alpha_3}\bigl(\underbrace{\alpha\tau_\infty^{\alpha_3}}_{\widetilde\alpha}\bigr)^{-1}
\tag{\unsure{*}}
\]
\mspage{244}

\scanpage{246}

where \deleted{$\unsure{\alpha\tau_\infty|\alpha_0}$}
\[
\widetilde\alpha=\alpha\tau_\infty^{\alpha_3}\deleted{\unsure{(\alpha_3)}}
\tag{45}
\]
has the advantage over $\alpha$ itself of being in $\widehat\Pi_{0,3}$ in all cases (\uline{N.B.} we have
$\alpha\in\widehat\Pi_{0,3}\cdot\{1,\tau_\infty\}$, and $\alpha\in\widehat\Pi_{0,3}$ iff $\alpha_3=0$).
It is equal to $\alpha$ iff $\alpha_3=0$ and in the case $\alpha_3=1$ \deleted{\illegible{1}} is directly related to $\xi$ by the formula
\[
\xi=\widetilde\alpha\ell_1^{-1}\rho(\widetilde\alpha)^{-1}
\tag{46}
\]
(deduced from (45), from $\xi=\alpha\rho(\alpha)^{-1}$, and from
\[
\tau_\infty\rho(\tau_\infty)^{-1}=\ell_1^{-1}
\tag{**}
\]
).
Thus (44) takes the definitive form
\[
\boxed{\widehat{\varphi}\bigl(\rho^{-1}(\xi)\bigr)=\rho^{j'}\sigma_\infty^{\alpha_3}\unsure{\lbrace}\widetilde\alpha^{-1}}
\qquad(\widetilde\alpha=\alpha\tau_\infty^{\alpha_3})
\tag{47}
\]
\margin{all the factors of this formula are \unsure{defined} in $\unsure{\widehat{\mathcal G}_{0,3}^{+}|\widetilde{\mathcal G}_{0,3}^{+}}$, i.e. of multiplier $+1$}
which can \unsure{also|thus} \unsure{be written|describe}, on passing to inverses
\[
\widehat{\varphi}\bigl(\rho^{-1}(\xi^{-1})\bigr)=\widetilde\alpha\sigma_\infty^{\alpha_3}\rho^{-j'}
\tag{47 bis}
\]
or again, solving for $\widetilde\alpha$
\[
\underset{\substack{\Vert\\\alpha\tau_\infty^{\alpha_3}}}{\widetilde\alpha}
=\widehat{\varphi}\bigl(\rho^{-1}(\xi^{-1})\bigr)\rho_{\deleted{\unsure{\infty}}}^{j'}\sigma_\infty^{\alpha_3}.
\tag{48}
\]
We shall compare this formula with formula (10), which is solved for $\beta$ in the form
\[
\widehat{\varphi}(\xi)=\beta\sigma_\infty^i
\tag*{(49)=(10)}
\]
\[
\beta=\widehat{\varphi}(\xi)\sigma_\infty^i
\tag{50}
\]
\mspage{245}

\scanpage{247}
\mspage{424}

The fact that formula (47) is indeed equivalent to (10), or again to (49), when $\alpha\in\unsure{\widehat{\widetilde\Pi}_{0,3}|\widetilde{\widehat\Pi}_{0,3}}$ and $\beta\in\widehat\Pi_{0,3}$ are related by $\xi=\eta\ell_1^\nu$ (where $\xi=\alpha\rho(\alpha)^{-1}$, $\eta=\beta\sigma_\infty(\unsure{\eta})^{-1}$), would deserve a direct verification. The interest of (47), or of \deleted{(48)} its form solved for $\alpha$ (48), just as of (10)\inserted{$=(49)$} or (50), consists in the fact that these formulas give the ``delicate'' parameters\inserted{$\alpha,\beta$}, in terms of the ``coarse parameter'' $\xi$. Both are formulas in $\unsure{\widehat{\mathcal G}_{0,3}^{+}}$ --- on reducing modulo $\widehat\Pi_{0,3}$, they give formulas in $\unsure{\mathfrak S_3|\mathcal S_3}$, namely (taking into account that $\widetilde\alpha,\beta\in\widehat\Pi_{0,3}$):
\deleted{\unsure{i.e.}}
\[
\left\{
\begin{aligned}
\widehat{\varphi}^{\cdot}\bigl(\rho^{-1}(\xi)\bigr)&=\dot\rho^{j'}\dot\sigma_\infty^{\alpha_3}\\
\widehat{\varphi}^{\cdot}(\xi)&=\dot\sigma_\infty^i
\end{aligned}
\right.
\tag{51}
\]
which \unsure{gives back|gives} exactly \deleted{\unsure{the}}, in terms of the elements $\widehat{\varphi}^{\cdot}(\xi)$ and $\widehat{\varphi}^{\cdot}\bigl(\rho^{-1}(\xi)\bigr)$ of $\unsure{\mathfrak S_3|\mathcal S_3}$, the values of $i,\alpha_3\in\mathbb Z/2\mathbb Z$ and of $j'\in\mathbb Z/3\mathbb Z$. (We should see whether $i,j'$ can, like $\alpha_3$, be made explicit in terms of the multiplier $\mu$, or even in terms of its
\margin{and also see whether $(i,\alpha_3,j')$ can take \unsure{all sorts} of values \unsure{in} $\mathbb Z/2\times\mathbb Z/2\times\mathbb Z/3$}
\mspage{246}

\scanpage{248}

quadratic residues $\varepsilon_\ell(\mu)$) We can supplement them with a formula deduced from one and \deleted{from} the other, taking into account that
\[
\xi\rho(\xi)\underbrace{\rho^2(\xi)}_{\rho^{-1}(\xi)}=1
\qquad\text{i.e.}\qquad
\begin{gathered}
\deleted{\rho^{-1}(\xi^{-1})=\xi\rho(\xi)}\\
\rho^{-1}(\xi)\xi=\rho(\xi)^{-1}
\end{gathered}
\]
we therefore find, multiplying (47) and (10)
\[
\widehat{\varphi}\bigl(\underbrace{\rho^{-1}(\xi)\cdot\xi}_{\rho(\xi)^{-1}}\bigr)
=\rho^{j'}\sigma_\infty^{\alpha_3}\widetilde\alpha^{-1}\beta\sigma_\infty^i
\]
i.e. again
\[
\widehat{\varphi}\bigl(\rho(\xi)\bigr)
=\sigma_\infty^i(\beta^{-1}\widetilde\alpha)\sigma_\infty^{\alpha_3}\rho^{-j'}
\tag{52}
\]
i.e., in the form solved for $\beta^{-1}\widetilde\alpha$
\[
\underbrace{\beta^{-1}\widetilde\alpha}_{(\beta^{-1}\alpha)\tau_\infty^{\alpha_3}}
=\sigma_\infty^i\widehat{\varphi}\bigl(\rho(\xi)\bigr)\sigma_\infty^{\alpha_3}\rho^{j'}.
\tag{53}
\]
We can consider that the three main ``delicate parameters'' are the quantities $\beta$, $\alpha$ and $\beta^{-1}\alpha$, which are expressed directly, in a remarkably simple way, in terms of the values of $\widehat{\varphi}$ for $\xi$, $\rho^{-1}(\xi)$ and $\unsure{\rho^{-1}|\rho^{-2}}(\xi)=\rho(\xi)$, which would dispense \deleted{\illegible{1}} probably (by
\mspage{247}

\scanpage{249}
\mspage{425}

carrying out the calculations in terms of $\xi,\eta$ exclusively, which would require a little more effort...) with appealing to delicate results \unsure{not} yet proved, \inserted{of the type $\xi=\alpha\rho(\alpha)^{-1}$, $\eta=\beta\sigma_\infty(\beta)^{-1}$,} to justify the introduction of these parameters $\alpha,\beta$.
\margin{(recapitulation)}
\[
\left\{
\begin{aligned}
\widehat{\varphi}(\xi)&=\beta\sigma_\infty^i\\
\widehat{\varphi}\bigl(\rho^{-1}(\xi)\bigr)&=\rho^{j'}\sigma_\infty^{\alpha_3}\widetilde\alpha^{-1}
\qquad(\widetilde\alpha=\alpha\tau_\infty^{\alpha_3})\\
\widehat{\varphi}\bigl(\underbrace{\rho^{\unsure{-2|-1}}(\xi)}_{\rho(\xi)}\bigr)
&=\underbrace{\sigma_\infty^i(\beta^{-1}\widetilde\alpha)\sigma_\infty^{\alpha_3}}_{\sigma_\infty^i(\beta^{-1}\widetilde\alpha)\sigma_\infty^{\alpha_3+i}\rho^{-j'}}\rho^{-j'}
\end{aligned}
\right.
\tag{54}
\]

Suppose for simplicity
\[
\alpha_3=0\quad\text{i.e.}\quad\alpha=\widetilde\alpha\quad\text{i.e.}\quad\alpha\in\widehat\Pi_{0,3}
\tag{55}
\]
Then we obtain from the second formula (54), taking inverses
\[
\widehat{\varphi}\bigl(\rho^{-1}(\xi^{-1})\bigr)=\alpha\rho^{-j'}
\]
and the $x\rho(x)^{-1}$ of both sides
\[
\xi=\widehat{\varphi}\bigl(\rho^{-1}(\xi^{-1})\bigr)\rho\Bigl(\widehat{\varphi}\bigl(\rho^{-1}(\xi^{-1})\bigr)\Bigr)^{-1}
\tag{56}
\]
which already implies the first formula (14) in the fundamental formulas (14) (15).
\mspage{248}

\scanpage{250}
\mspage{425}

We thus find in this case ($\mu\equiv1\ (3)$ i.e.
$\nu\equiv0\ (3)$) that the two relations (14) (15)
are also equivalent to \unsure{the set|one} of the two relations
\margin{Valid only if $\nu\equiv0\ (3)$ i.e.
$\mu\equiv1\ (3)$ i.e. $\varepsilon_3(\mu)=1$.}
\[
\begin{array}{c@{\quad}c}
\begin{gathered}(57)\\[1.5ex](58)\end{gathered}&
\boxed{\begin{aligned}
\xi&=
\widehat\varphi\bigl(\rho^{-1}(\xi^{-1})\bigr)
\rho\bigl(\widehat\varphi(\rho^{-1}(\xi^{-1}))\bigr)^{-1}
\\
\xi&=\widehat\varphi(\xi)\,
\sigma_\infty\bigl(\widehat\varphi(\xi)\bigr)^{-1}
\ell_1^\nu
\end{aligned}}
\end{array}
\]
\inserted{with $\nu\in\widehat{\mathbb Z}$,
$\nu\equiv0\ (3)$}
which therefore characterize the
$\unsure{u_0|u_{\cdot}}=u_{\xi}\in
\unsure{\mathcal M_{0,3}^{!\sim}|\mathcal M_{0,3}^{!\approx}}(0)$,
whose multiplier $\mu$ is $\equiv1\ (3)$.

\uline{Remark} The only element of
$\mathcal M_{0,3}^{!\approx}(0)$
\inserted{$\simeq\mathcal N_{0,3}^{!\approx}$}
that I know how to make explicit, apart from the unit element
(and \unsure{one} even in
$\mathcal N_{0,3}^{!\sim}\supset
\unsure{\mathcal N_{4,3}^{!\approx}|\mathcal N_{g,3}^{!\approx}}$)
is \deleted{\illegible{1}}\inserted{the unique} element
$\tau\in\mathcal T_{0,3}^{!}$
\inserted{$\subset\mathcal N_{0,3}^{!\approx}$},
which in $\mathcal M_{0,3}^{!\approx}(0)$ corresponds
\inserted{to the element $\tau_\infty$, \deleted{\illegible{1}}
$=u_{\xi}=u_{\xi,\eta}$}
\deleted{to the invariants}
\[
\begin{aligned}
(59)\qquad & h=\beta_0=1,\quad f=1,\quad
\alpha=\alpha_0=\tau_\infty,\quad
\xi=g=\ell_1^{-1},\quad\eta=1
\\
&\nu=-1,\quad\mu=-1,
\end{aligned}
\]
\margin{Careful, here $\alpha_3=1$ so (57) is not
\unsure{valid|applicable}.}
\deleted{which}\inserted{\unsure{and} $\xi,\nu$}
do satisfy \deleted{\unsure{(57) (58)}}
\uline{$\xi\rho(\xi)
\rho^2(\xi)=1$}
\inserted{i.e. $\ell_0^{-1}\ell_1^{-1}\ell_\infty^{-1}=1$}
\margin{and the relations (54) \unsure{page above}}
\deleted{
\[
\widehat\varphi(\xi)=\ell_0,\qquad
\widehat\varphi\bigl(\unsure{\rho^{-1}}(\xi^{-1})\bigr)
=\unsure{\varepsilon_0},
\]
\[
\widehat\varphi\bigl(\rho^{-1}(\xi^{-1})\bigr)
\rho\bigl(\widehat\varphi(\rho^{-1}(\xi^{-1}))\bigr)^{-1}
=\unsure{\varepsilon_0\varepsilon_1^{-1}}
=(\sigma_\infty\rho)(\rho\sigma_\infty)^{-1}.
\]
}
\mspage{249}

\scanpage{251}
\mspage{426}

with
\[
(60)\qquad i=1,\quad\alpha_3=1\quad\text{i.e.}\quad
\varepsilon_3=-1,\quad j'=-1
\]

\uline{Remark} We can ask the question, in general,
of determining the elements
$u_{\alpha,\beta}=u_{\xi,\unsure{\eta}}$
of $\mathcal M_{0,3}^{!\sim}[0]$
(not necessarily normalized in one sense or the other)
such that we have (of course by hypothesis, $f=1$ ; but also)
\deleted{either} $\beta=1$ i.e. $\eta=1$,
\deleted{or} \deleted{\illegible{3}}
(\uline{N.B.} We already know that if $\alpha=1$ i.e.
$\xi=1$, then $u=1$)
We shall have to have
$\xi=\eta\ell_1^\nu=\ell_1^\nu$,
and it remains to check
\deleted{\illegible{2}}
\[
\xi\rho(\xi)
\rho^2(\xi)=1\quad\text{with}\quad
\xi=\ell_1^\nu,
\]
$\unsure{\mathrm{Sl}(2,\widehat{\mathbb Z})'}$,
\deleted{this} so that
\[
\underline{\xi}=
\begin{pmatrix}1&0\\2\nu&1\end{pmatrix}
\]
this means \inserted{(page \unsure{344|342})} that we have either
\[
\underline{\xi}=\pm\underline\rho^{-1}.
\quad\deleted{\text{i.e.}\ \underline{\xi}=\rho}
=\pm\begin{pmatrix}-1&1\\1&0\end{pmatrix}
\]
-- which is clearly not the case -- or
\[
C+D-B=\pm1\quad\text{i.e.}\quad2\nu+1=\mu=\pm1\quad\text{i.e.}
\]
$\nu=0$ or $-1$ -- these are the two trivial cases
$u_{\xi}=\mathrm{id}$ and
$u_{\xi}=\tau_\infty$.

\margin{We are \unsure{needlessly} repeating arguments already
made, we know very well that the multiplier is
(up to sign when the \unsure{normalization} of
$\xi$ has not been \unsure{done}) equal to one
of the ``\unsure{coefficients} \illegible{1}'' of
$\xi$, hence equal to $\pm1$.}
\mspage{250}

\scanpage{252}

Let us note \unsure{the|the}

\uline{Proposition.} The commutant of $\tau$ (the
nontrivial element of $\mathcal T_{0,3}^{!}\simeq\mathbb Z/2\mathbb Z
\subset\unsure{\mathcal N_{0,3}^{!\sim}|\mathcal M_{0,3}^{!\sim}}$)
in $\mathcal N_{0,3}^{!\sim}$ is reduced to $\{1,\tau\}$.

\margin{\uline{N.B.} This proposition \illegible{3}
in $\S\,46$ \unsure{VI|II}, \illegible{4} of \illegible{2}
of $\unsure{\widehat\Gamma}$.}

We shall do the calculation in the isomorphic group
$\unsure{\mathcal M_{0,3}^{!\sim}}(0)$,
\inserted{where $\tau$ corresponds to $\tau_\infty$},
in fact it turns out that the centralizer of $\tau$ in
the whole of $\mathcal M_{0,3}^{!\sim}[0]$
(the big \unsure{normalizer} of $L_0$ in
$\mathcal M_{0,3}^{!\sim}$, formed of the
$u_{\xi}$ without normalization conditions)
is reduced to $\{1,\tau\}$.
Indeed, we must express the commutation of a
$u=u_{\xi}$ with
$u_{\xi'}=\tau_\infty$,
where $\xi'=\ell_1^{-1}$, we have
\[
\left\{\begin{aligned}
u_{\xi'}u_{\xi}
&=u_{\xi''}
&&\text{with }\xi''=
\dot u_{\xi'}(\xi)\xi'
\\
u_{\xi}u_{\xi'}
&=u_{\xi'''}
&&\text{with }\xi'''=
\unsure{\dot u|u}_{\xi}(\xi')\xi
\end{aligned}\right.
\]
and we must express $\xi''=\xi'''$, i.e.
\[
(61)\qquad
\tau_\infty(\xi)\ell_1^{-1}
=u_{\xi}(\ell_1^{-1})\xi.
\]
Now we have
\[
\unsure{\dot u|u}_{\xi}(\varepsilon_1)
=\mathrm{int}(\xi)(\varepsilon_1^\mu)
\qquad\text{whence}
\]
\[
\unsure{\dot u|u}_{\xi}
\bigl(\underbrace{\ell_1^{-1}}_{\varepsilon_1^{-2}}\bigr)
=\mathrm{int}(\xi)(\ell_1^{-\mu})
\]
and (61) is written
\[
\tau_\infty(\xi)\ell_1^{-1}
=\xi\ell_1^{-\mu}\qquad\text{i.e.}
\]
\[
(62)\qquad
\xi^{-1}\tau_\infty(\xi)
=\ell_1^{1-\mu}=\ell_1^{-2\nu}.
\]
\mspage{251}